%!TEX root = main.tex 

\section{Limitations and Discussions} 
\label{sec:discussion}

\noindent\sys has limitations and we wish to improve it in at least the following directions.
\begin{comment}
\noindent\textbf{Computation prediction.}
\sys relies on trace-based profiling rather than purely architectural models. By extracting real execution traces from small-scale runs and extrapolating to large systems, \sys avoids the high inaccuracies (e.g., >70\%) we observed when using architecture-focused simulators like SCALE-sim~\cite{samajdar2018scale}. In the future, we plan to explore hybrid approaches that blend \sys's profiling accuracy with more flexible architectural models.
\end{comment}

\noindent\textbf{Communication prediction.}
% \sys's white-box modeling and offline profiling offer an accurate, low-overhead method for communication kernel estimation. 
\sys positions itself between hand-tuned $\alpha$--$\beta$ models and packet-level network simulators.
We plan to investigate learning based solutions like m3~\cite{li2024m3} that efficiently predicts flow-level performance while capturing network-centric factors such as transport protocols. 
% We also wish to extend our modeling to explicitly consider network topology and the potential bandwidth contention, which is currently missing.

\noindent\textbf{Parallelism.}
\sys does not support sequence/context parallelism currently.
We plan to integrate them in subsequent releases.

\noindent\textbf{MoE routing policies.}
Our current paired validation focuses on the widely used no-token-dropping path in Megatron-style MoE training and evaluates controlled routing skew as the primary source of expert-load imbalance. 
We leave explicit capacity-factor extremes, token-dropping/capacity-aware admission control, and their effects on step-time tails and straggler behavior to future work.


% \yc{TODO: We also provide the comparison with packet-level simulator SimAI for further discussion of trad-off between fidelity and simulation speed.}



% It is also very interesting to 


% communication predictor currently relies on the knowledge profiled on a four-server cluster. Conventionally, it is the minimum requirement to construct the cluster suitable for double binary tree-based all-reduce. It is feasible to lower the required number of servers and use two servers to imitate such cluster topology. However, since the knowledge can be profiled ahead-of-time and is mostly reusable, it does not affect the results of the simulation. We leave this potential improvement to future work so that the profiling cost can be further reduced.


% \noindent\textbf{Model parallelism.} In this work, we discuss the simulation of data-parallel distributed DNN, which is the primary practice in DNN training. Model parallelism~\cite{shoeybi2019megatron, huang2019gpipe, NHPS19}is adopted for DNN models that cannot fit into one GPU memory. Model-parallel distributed training does not change the problem nature but requires additional send and receive operations among devices to transmit intermediate output. Through pipelining, they are also overlapped with computation kernels. Our prediction model can be trained to learn the interference in this scenario as model parallelism does not change the problem nature.

